\chapter{EXPERIMENTAL PROTOTYPE}
\label{ch:prototype}

\section{Introduction}

The previous chapters described what \projectname{} is meant to do and why the
market has room for it. This chapter describes what was actually built. It
covers the environment and the technologies used, the architecture of the
implemented system, the internals of the matching engine --- which is the part
of the platform that carries the innovation claimed in
Chapter~\ref{ch:innovation} --- the user interfaces, and the way the whole thing
is tested and packaged for deployment.

Two decisions shape everything that follows and are worth stating up front.
First, \textbf{no model inference ever happens inside an HTTP request}: encoding
a photograph with a vision transformer takes long enough that doing it while the
user waits would make reporting an item feel broken, so all inference is pushed
onto a background worker and the API answers immediately. Second,
\textbf{vectors live in the main database}, not in a separate vector store: the
core query of this platform is \dq{find items of the opposite type, still open,
reported within a plausible date window, whose description or photograph
resembles this one}, and that is one query joining relational predicates to a
similarity search. A dedicated vector index such as FAISS~\cite{johnson2021faiss}
or a managed vector database would be faster in isolation, but it cannot apply
those relational predicates, so the system would have to retrieve loosely and
filter afterwards in application code --- discarding most of what it fetched and
losing transactional consistency between an item and its vector. At the scale
this platform will reach in its first years, that trade is not worth making.

\section{Development environment, tools and technologies}

\subsection{Hardware environment and editors}

Development was carried out on a personal workstation running Windows~11, with
the entire backend stack executing inside Linux containers through Docker
Desktop. Model inference runs on CPU: both encoders used by the platform are
small enough that a GPU is a convenience rather than a requirement, and building
for CPU keeps the deployment target realistic for the low-cost hosting described
in Chapter~\ref{ch:finance}. Source editing was done in Visual Studio Code, with
the Python and TypeScript language servers providing type checking against the
same configurations used in continuous integration.

\subsection{Collaborative tools and version control}

The project is a single Git repository containing the backend, the frontend, the
infrastructure definitions and the design documentation, so that a change
touching an API contract and its client can be reviewed as one unit. GitHub
hosts the remote. The API documentation is not written by hand: FastAPI derives
an OpenAPI schema from the Pydantic models, and an interactive Swagger interface
is served alongside the running API, which means the contract shown to the
frontend developer is generated from the code that enforces it and cannot drift
away from it.

\subsection{Client-side technologies (front-end)}

The interface is built with Next.js~\cite{nextjs} and styled with Tailwind
CSS~\cite{tailwind}; Table~\ref{tab:frontend-stack} lists the full stack and the
role each part plays.

\begin{table}[htbp]
\centering
\caption{Front-end technology stack}
\label{tab:frontend-stack}
\small
\begin{tabularx}{\textwidth}{@{}l l X@{}}
\toprule
\textbf{Concern} & \textbf{Technology} & \textbf{Role in the prototype} \\
\midrule
Framework    & Next.js 15 (App Router) & Server and client rendering, file-system routing, locale segment \code{/[locale]} \\
UI library   & React 19 + TypeScript   & Component model with end-to-end static typing against the API types \\
Styling      & Tailwind CSS 3.4        & Utility-first styling; logical properties give right-to-left support at no extra cost \\
Components   & shadcn/ui on Radix UI   & Accessible primitives (dialog, select, dropdown, switch, avatar) owned in-repo rather than imported as a black box \\
Server state & TanStack Query 5        & Caching, revalidation and optimistic updates for every API read \\
Client state & Zustand                 & Session and UI state that is not server-owned \\
Forms        & React Hook Form + Zod   & Declarative validation; the Zod schema is the single source of truth for the report form \\
i18n         & next-intl               & Message catalogues for English, French and Arabic, with locale-prefixed routes \\
Theming      & next-themes             & Light and dark presentation \\
Motion       & Framer Motion           & Transitions on the report flow and match reveal \\
Feedback     & Sonner                  & Toast notifications \\
\bottomrule
\end{tabularx}
\end{table}

\subsection{Server-side technologies (back-end)}

The API is written with FastAPI~\cite{fastapi} over SQLAlchemy~\cite{sqlalchemy}
and PostgreSQL~16~\cite{postgresql}, with Redis~\cite{redis} carrying the
\code{arq}~\cite{arq} job queue and Docker~\cite{docker} packaging the whole
system. Authentication, upload handling and rate limiting follow the OWASP Top
Ten guidance~\cite{owasp2021}.

\begin{table}[htbp]
\centering
\caption{Back-end technology stack}
\label{tab:backend-stack}
\small
\begin{tabularx}{\textwidth}{@{}l l X@{}}
\toprule
\textbf{Concern} & \textbf{Technology} & \textbf{Role in the prototype} \\
\midrule
API framework & FastAPI          & Asynchronous HTTP layer, dependency injection, generated OpenAPI schema \\
Validation    & Pydantic v2      & Request and response models; settings loaded and type-checked from the environment \\
ORM           & SQLAlchemy 2 (async) + asyncpg & Declarative models with typed mappings, executed on an async driver \\
Migrations    & Alembic          & Seven versioned migrations, including the \code{pgvector} column and index definitions \\
Database      & PostgreSQL 16 + \code{pgvector} & Relational data, full-text search and vector similarity in one engine \\
Queue         & Redis 7 + \code{arq} & Job queue for embedding and matching; \code{arq} is asyncio-native, matching the API's concurrency model \\
Auth          & PyJWT + Argon2   & Short-lived access tokens, rotating refresh tokens, memory-hard password hashing \\
Images        & Pillow           & Decoding, validation and re-encoding of uploads to WebP \\
Storage       & \code{StorageBackend} abstraction & Local filesystem in development, S3-compatible object storage in production, behind one interface \\
\bottomrule
\end{tabularx}
\end{table}

\subsection{Machine-learning stack}
\label{sec:ml-stack}

Two pre-trained encoders are used, both loaded through the
\code{sentence-transformers} library~\cite{sentencetransformers} so that a
single API (\code{.encode()}) covers text and images alike, and both pinned to
CPU builds of PyTorch.

\paragraph{Text --- \code{paraphrase-multilingual-MiniLM-L12-v2} (384 dimensions).}
This is a Sentence-BERT model~\cite{reimers2019sbert} --- a
BERT~\cite{devlin2019bert} transformer~\cite{vaswani2017attention} fine-tuned
so that whole sentences, rather than individual tokens, occupy a space where
cosine distance is meaningful --- made multilingual by
knowledge distillation from an English teacher onto a multilingual student, so
that translations of the same sentence are mapped to nearby
vectors~\cite{reimers2020multilingual}.
The choice of a \emph{multilingual} sentence encoder is not incidental; it is
forced by the market described in Chapter~\ref{ch:market}. Reports on this
platform arrive in Arabic, French and English, and frequently mix two of them
inside a single description. The English-only encoder that is usually quoted for
this kind of task, \code{all-MiniLM-L6-v2}, places \dq{Téléphone Samsung égaré}
and \dq{Samsung phone lost} almost as far apart as two unrelated sentences ---
which is precisely the match the product exists to find. The multilingual model
covers more than fifty languages and happens to produce vectors of the same 384
dimensions, so it substitutes for the English model without any change to the
database schema.

\paragraph{Images --- CLIP \code{ViT-B/32} (512 dimensions).}
CLIP~\cite{radford2021clip} is a vision transformer~\cite{dosovitskiy2021vit}
trained contrastively on image--caption pairs.
Photographs are compared only to other photographs, so no language is involved
on the image side and the standard CLIP checkpoint is the correct choice. CLIP
is preferred to the penultimate layer of a supervised classifier because its
representation is general-purpose: it was never trained on a fixed label set,
and lost property is an open-ended category by definition.

\paragraph{Why not one model for both.} CLIP embeds images and text into a
shared space, which is attractive at first glance. But the dominant comparison
in this system is text against text --- one person's description of a wallet
against another person's description of the same wallet --- and for that a
dedicated sentence encoder is both stronger and cheaper than CLIP's text tower.
The two-model arrangement uses each where it is strongest. CLIP's text encoder
remains available for a future cross-modal signal, matching the \emph{words} of
one report against the \emph{photograph} of another.

\section{Architecture of the implemented system}

\subsection{Overall architecture}

The backend is a \textbf{modular monolith}: one FastAPI application whose
internal boundaries (\code{api/}, \code{services/}, \code{repositories/},
\code{ml/}, \code{storage/}) are drawn so that any module could later be
extracted into its own service, plus a background worker that shares the same
codebase and image but runs a different entry point.
Figure~\ref{fig:architecture} shows the deployed topology.

\begin{figure}[htbp]
\centering
\begin{tikzpicture}[x=1cm,y=1cm]

  \node[boxnode, text width=10.4cm] (client) at (0,6.2)
    {\textbf{Next.js 15 --- browser and server rendering}\\
     locale-prefixed routes \dcode{/en} \dcode{/fr} \dcode{/ar} $\cdot$ TanStack Query cache};

  \node[boxnode, text width=10.4cm] (api) at (0,3.9)
    {\textbf{FastAPI --- \dcode{/api/v1}}\\
     auth $\cdot$ categories $\cdot$ items $\cdot$ images $\cdot$ matches $\cdot$ claims $\cdot$ notifications\\
     authentication $\cdot$ validation $\cdot$ authorisation $\cdot$ OpenAPI};

  \node[storenode, text width=2.7cm] (redis) at (-4,1.3)
    {\textbf{Redis 7}\\ job queue, cache};
  \node[storenode, text width=2.7cm] (pg) at (0,1.3)
    {\textbf{PostgreSQL 16}\\ \dcode{pgvector} + FTS};
  \node[storenode, text width=2.7cm] (store) at (4,1.3)
    {\textbf{Object storage}\\ local FS $\rightarrow$ S3};

  \node[boxnode, text width=10.4cm] (worker) at (0,-1.2)
    {\textbf{\dcode{arq} worker --- all model inference}\\
     \dcode{embed\_item} $\rightarrow$ \dcode{run\_matching} $\rightarrow$ notifications\\
     multilingual MiniLM (384-d) $\cdot$ CLIP ViT-B/32 (512-d)};

  \draw[flowarrow] (client) -- node[edgelabel,right]{HTTPS, JSON, JWT bearer} (api);

  \draw[flowarrow] (api.south -| redis.north) -- (redis.north)
    node[edgelabel,midway,left]{enqueue};
  \draw[flowarrow] (api.south -| pg.north) -- (pg.north)
    node[edgelabel,midway,right]{read / write};
  \draw[flowarrow] (api.south -| store.north) -- (store.north)
    node[edgelabel,midway,right]{put / get};

  \draw[flowarrow] (redis.south) -- (worker.north -| redis.south)
    node[edgelabel,midway,left]{dequeue};
  \draw[flowarrow] (worker.north -| pg.south) -- (pg.south)
    node[edgelabel,midway,right]{write vectors, matches};
  \draw[flowarrow] (worker.north -| store.south) -- (store.south)
    node[edgelabel,midway,right]{read images};

\end{tikzpicture}
\caption{Deployed architecture. The API never loads a model; the worker never
serves a request.}
\label{fig:architecture}
\end{figure}

\subsection{Why inference runs outside the request cycle}

Encoding one photograph with CLIP takes on the order of hundreds of
milliseconds on CPU, and the encoders together occupy well over a gigabyte of
resident memory. Performing that work inside \code{POST /items} would mean three
separate costs: the user waits for a neural network before seeing any
confirmation, the API process must hold the model weights, and API latency
becomes coupled to model load time on every cold start.

Separating the worker removes all three. The API validates the report, persists
it, returns \code{201~Created}, and enqueues a job. The item is publicly
browsable immediately with \code{processing\_status~=~pending}, and the worker
advances it through \code{embedding} and \code{matching} to \code{ready},
setting \code{failed} on any unhandled error so that a sweep can retry it. The
business lifecycle of the item (\code{open}, \code{matched}, \code{claimed},
\code{closed}) is deliberately kept as a separate column from this pipeline
state, because they answer different questions: an item can be perfectly
browsable while its matches are still being computed.

\subsection{Implemented data model}

Lost items and found items are stored in \textbf{one \code{items} table
discriminated by a \code{type} enum}, not in two tables. They share every column,
the same embedding pipeline and the same lifecycle; the only difference is the
direction in which matching searches. Two tables would have duplicated the
schema, the indexes and the embedding code, and forced a \code{UNION} into every
browse query. The conceptual entities survive as filtered views.

\begin{table}[htbp]
\centering
\caption{Implemented entities}
\label{tab:entities}
\small
\begin{tabularx}{\textwidth}{@{}l X@{}}
\toprule
\textbf{Table} & \textbf{Purpose and notable columns} \\
\midrule
\code{users} & Account, role (\code{user}/\code{admin}), verification and active flags. Email stored as \code{citext} so that case never creates a duplicate account. \\
\code{refresh\_tokens} & Rotation and revocation. Only the \emph{hash} of a token is stored; \code{replaced\_by} forms a chain used for reuse detection. \\
\code{categories} & Seeded taxonomy (phones, wallets, keys, bags, documents, \ldots), self-referencing for subcategories. Doubles as a hard pre-filter during matching. \\
\code{items} & Both report types. \code{type}, \code{title}, \code{description}, \code{category\_id}, \code{color}, \code{brand}, \code{wilaya\_code}, \code{lost\_or\_found\_at}, business \code{status}, ML \code{processing\_status}, and \code{text\_embedding vector(384)}. \\
\code{item\_images} & Up to five photographs per item; \code{image\_path} identifies the stored file, \code{image\_embedding vector(512)} holds the CLIP vector, and \code{blur\_preview} supports locked match previews. \\
\code{matches} & One row per (lost, found) pair above the persistence threshold: \code{text\_score}, \code{image\_score}, \code{combined\_score}, \code{confidence}, \code{explanation} (JSONB), \code{status}. Unique on the pair, so re-matching upserts instead of duplicating. \\
\code{match\_feedback} & One user's correctness verdict for a match, retained as labelled data for matching calibration. \\
\code{claims} & The ownership handshake: \code{answers} (JSONB, questions snapshotted at claim time), \code{status}, \code{resolved\_at}. A partial unique index permits only one \emph{pending} claim per person per item. \\
\code{notifications} & In-app messages with an optional enforced link to an item and an optional match UUID for navigation. \\
\code{payments} & A payment-gateway checkout attempt, including the price and credit-pack snapshot, provider reference, status, and provider payload. \\
\code{credit\_ledger} & Append-only credit movements. The balance is \code{SUM(delta)}, rather than a mutable column on \code{users}. \\
\code{match\_unlocks} & Records that a user has permanently unlocked a match; the user--match pair is unique. \\
\bottomrule
\end{tabularx}
\end{table}

Figure~\ref{fig:erd} shows every table in the implemented schema and its
database-enforced foreign-key relationships.

\begin{sidewaysfigure}[p]
\centering
\begin{tikzpicture}[x=1cm,y=1cm]

  % Account and taxonomy tables.
  \node[weakentity] (rt) at (-7.8,6.3)
    {\textbf{refresh\_tokens}\\ \dcode{id} PK\\ \dcode{user\_id} FK\\
     \dcode{token\_hash} UK\\ \dcode{expires\_at}\\ \dcode{revoked\_at}\\
     \dcode{replaced\_by}\\ \dcode{user\_agent}\\ \dcode{ip}\\
     \dcode{created\_at}/\dcode{updated\_at}};

  \node[entity] (usr) at (-7.8,2.5)
    {\textbf{users}\\ \dcode{id} PK\\ \dcode{email} UK\\ \dcode{password\_hash}\\
     \dcode{full\_name}\\ \dcode{phone}\\ \dcode{role}\\ \dcode{status}\\
     \dcode{avatar\_url}\\ \dcode{is\_verified}\\
     \dcode{created\_at}/\dcode{updated\_at}};

  \node[entity] (cat) at (-7.8,-1.3)
    {\textbf{categories}\\ \dcode{id} PK\\ \dcode{parent\_id} FK\\
     \dcode{name}\\ \dcode{slug} UK\\ \dcode{created\_at}/\dcode{updated\_at}};

  \node[entity] (nt) at (-7.8,-4.7)
    {\textbf{notifications}\\ \dcode{id} PK\\ \dcode{user\_id} FK\\
     \dcode{type}\\ \dcode{title}\\ \dcode{body}\\ \dcode{is\_read}\\
     \dcode{item\_id} FK\\ \dcode{match\_id}\\
     \dcode{created\_at}/\dcode{updated\_at}};

  % Reports and their dependent records.
  \node[entity] (img) at (-2.6,6.5)
    {\textbf{item\_images}\\ \dcode{id} PK\\ \dcode{item\_id} FK\\
     \dcode{image\_path}\\ \dcode{image\_embedding} \textit{vector(512)}\\
     \dcode{blur\_preview}\\ \dcode{created\_at}/\dcode{updated\_at}};

  \node[entity] (itm) at (-2.6,1.8)
    {\textbf{items}\\ \dcode{id} PK\\ \dcode{user\_id} FK\\ \dcode{category\_id} FK\\
     \dcode{type} \textit{lost|found}\\ \dcode{status}\\ \dcode{processing\_status}\\
     \dcode{title}\\ \dcode{description}\\ \dcode{search\_vector}\\
     \dcode{color}\\ \dcode{brand}\\ \dcode{location\_text}\\
     \dcode{wilaya\_code}\\ \dcode{latitude}/\dcode{longitude}\\
     \dcode{lost\_or\_found\_at}\\ \dcode{claim\_questions} \textit{jsonb}\\
     \dcode{closed\_reason}/\dcode{closed\_at}\\
     \dcode{text\_embedding} \textit{vector(384)}\\
     \dcode{created\_at}/\dcode{updated\_at}};

  \node[entity] (cl) at (-2.6,-4.5)
    {\textbf{claims}\\ \dcode{id} PK\\ \dcode{item\_id} FK\\
     \dcode{claimant\_id} FK\\ \dcode{status}\\ \dcode{message}\\
     \dcode{answers} \textit{jsonb}\\ \dcode{resolved\_at}\\
     \dcode{created\_at}/\dcode{updated\_at}};

  % Matching and payment records.
  \node[entity] (mt) at (2.6,5.7)
    {\textbf{matches}\\ \dcode{id} PK\\ \dcode{lost\_item\_id} FK\\
     \dcode{found\_item\_id} FK\\ \dcode{status}\\ \dcode{text\_score}\\
     \dcode{image\_score}\\ \dcode{combined\_score}\\ \dcode{confidence}\\
     \dcode{explanation} \textit{jsonb}\\ \dcode{resolved\_at}\\
     \dcode{created\_at}/\dcode{updated\_at}};

  \node[weakentity] (mf) at (2.6,1.7)
    {\textbf{match\_feedback}\\ \dcode{id} PK\\ \dcode{match\_id} FK\\
     \dcode{user\_id} FK\\ \dcode{is\_correct}\\ \dcode{comment}\\
     \dcode{created\_at}/\dcode{updated\_at}};

  \node[entity] (pay) at (2.6,-3.1)
    {\textbf{payments}\\ \dcode{id} PK\\ \dcode{user\_id} FK\\ \dcode{pack\_id}\\
     \dcode{credits}\\ \dcode{amount}\\ \dcode{currency}\\ \dcode{status}\\
     \dcode{provider}\\ \dcode{provider\_ref}\\ \dcode{checkout\_url}\\
     \dcode{provider\_payload} \textit{jsonb}\\ \dcode{failure\_reason}\\
     \dcode{paid\_at}\\ \dcode{created\_at}/\dcode{updated\_at}};

  \node[weakentity] (led) at (7.8,4.1)
    {\textbf{credit\_ledger}\\ \dcode{id} PK\\ \dcode{user\_id} FK\\
     \dcode{payment\_id} FK\\ \dcode{match\_id} FK\\ \dcode{delta}\\
     \dcode{reason}\\ \dcode{note}\\ \dcode{created\_at}/\dcode{updated\_at}};

  \node[weakentity] (unl) at (7.8,0.4)
    {\textbf{match\_unlocks}\\ \dcode{id} PK\\ \dcode{user\_id} FK\\
     \dcode{match\_id} FK\\ \dcode{source}\\
     \dcode{created\_at}/\dcode{updated\_at}};

  % Each line is a real foreign key. Lines are kept behind nodes so crossings
  % never make any table field difficult to read.
  \begin{scope}[on background layer]
    \draw[relline] (usr.north) -- (rt.south)
      node[card,pos=0.5,right] {1 : N};
    \draw[relline] (usr.east) -- (itm.west)
      node[card,pos=0.52,above] {1 : N};
    \draw[relline] (cat.east) -- (itm.west)
      node[card,pos=0.55,below] {1 : N};
    \draw[relline] (itm.north) -- (img.south)
      node[card,pos=0.5,right] {1 : N};
    \draw[relline] (usr.south) -- (nt.north)
      node[card,pos=0.5,right] {1 : N};
    \draw[relline] (itm.south west) to[out=-150,in=30]
      (nt.east) node[card,pos=0.5,below] {1 : N};
    \draw[relline] (itm.south) -- (cl.north)
      node[card,pos=0.5,right] {1 : N};
    \draw[relline] (usr.south east) to[out=-25,in=155]
      (cl.west) node[card,pos=0.5,below] {1 : N};
    \draw[relline] (itm.north east) to[out=25,in=180]
      (mt.west) node[card,pos=0.52,above] {1 : N (lost)};
    \draw[relline] (itm.east) to[out=5,in=210]
      (mt.south west) node[card,pos=0.52,below] {1 : N (found)};
    \draw[relline] (mt.south) -- (mf.north)
      node[card,pos=0.5,right] {1 : N};
    \draw[relline] (usr.east) to[out=18,in=180]
      (mf.west) node[card,pos=0.5,above] {1 : N};
    \draw[relline] (usr.south east) to[out=-18,in=165]
      (pay.west) node[card,pos=0.5,below] {1 : N};
    \draw[relline] (usr.east) to[out=8,in=180]
      (led.west) node[card,pos=0.55,above] {1 : N};
    \draw[relline] (pay.east) -- (led.south west)
      node[card,pos=0.5,below] {1 : N};
    \draw[relline] (mt.east) -- (led.west)
      node[card,pos=0.5,above] {1 : N};
    \draw[relline] (usr.south east) to[out=-12,in=180]
      (unl.west) node[card,pos=0.58,below] {1 : N};
    \draw[relline] (mt.south east) to[out=-20,in=130]
      (unl.north west) node[card,pos=0.55,right] {1 : N};
    % A category may be the parent of many child categories.
    \draw[relline] (cat.west) .. controls +(-1.2,0) and +(-1.2,-1.0) ..
      (cat.south) node[card,pos=0.5] {1 : N parent};
  \end{scope}
\end{tikzpicture}
\caption{Complete entity-relationship diagram of the implemented schema.
Every line represents an enforced foreign key; \code{notifications.match\_id}
is deliberately not connected because it is stored as an unconstrained UUID,
not a foreign key. \code{items} participates twice in \code{matches}: once as
the lost report and once as the found report.}
\label{fig:erd}
\end{sidewaysfigure}

Migrations are managed with Alembic across eight versions. Migrations 0005--0007
add full-text search plus the text and image vector indexes; migration 0008 adds
the payment, credit-ledger, and match-unlock tables. Both vector indexes are
HNSW~\cite{malkov2020hnsw} over the cosine operator
class provided by \code{pgvector}~\cite{pgvector}, which gives better
recall-to-latency behaviour than IVFFlat at the corpus sizes this platform will
see in its first years.

\section{The matching engine}

This section is the technical core of the prototype. The engine answers one
question: \emph{given a newly reported item, which items of the opposite type
are most likely to be the same physical object?} Figure~\ref{fig:pipeline}
summarises the pipeline.

\begin{figure}[htbp]
\centering
\begin{tikzpicture}[x=1cm,y=1cm]

  \node[boxnode, text width=6cm] (new) at (0,7.2)
    {\textbf{New or edited report} (lost or found)};

  \node[boxnode, text width=4.4cm] (txt) at (-3.1,5.4)
    {Multilingual MiniLM\\ title + description + metadata $\rightarrow$ 384-d};
  \node[boxnode, text width=4.4cm] (img) at (3.1,5.4)
    {CLIP ViT-B/32\\ each photograph $\rightarrow$ 512-d};

  \node[storenode, text width=8.6cm] (store) at (0,3.6)
    {Vectors written to PostgreSQL (\dcode{pgvector}, unit-normalised)};

  \node[boxnode, text width=8.6cm] (retr) at (0,1.8)
    {\textbf{Retrieval} --- relational pre-filter, then two independent routes:\\
     text ANN {\footnotesize(+ lexical FTS)} $\cup$ image ANN};

  \node[boxnode, text width=8.6cm] (fuse) at (0,0)
    {\textbf{Scoring} --- lexical blend, image calibration, weighted fusion};

  \node[boxnode, text width=8.6cm] (conf) at (0,-1.8)
    {\textbf{Confidence} --- metadata boosts, distinctiveness margin, thresholds};

  \node[boxnode, text width=8.6cm] (out) at (0,-3.6)
    {Upsert \dcode{matches}, retract stale ones, notify both owners};

  \draw[flowarrow] (new.south) -- ++(0,-0.35) -| (txt.north);
  \draw[flowarrow] (new.south) -- ++(0,-0.35) -| (img.north);
  \draw[flowarrow] (txt.south) |- ($(store.north)+(0,0.35)$) -- (store.north);
  \draw[flowarrow] (img.south) |- ($(store.north)+(0,0.35)$) -- (store.north);
  \draw[flowarrow] (store) -- (retr);
  \draw[flowarrow] (retr) -- (fuse);
  \draw[flowarrow] (fuse) -- (conf);
  \draw[flowarrow] (conf) -- (out);

\end{tikzpicture}
\caption{The matching pipeline, from a new report to persisted suggestions.}
\label{fig:pipeline}
\end{figure}

\subsection{Generating the embeddings}

The worker assembles one document string from the structured and free-text
fields of a report and encodes it in a single pass. Normalising at encode time
makes every vector unit length, so that the cosine similarity between two
vectors equals their dot product and \code{pgvector}'s cosine distance operator
is directly usable without further scaling.

\begin{lstlisting}[language=Python,caption={Building and encoding the text document},label={lst:embed}]
doc = f"{title}. {description}. category: {category_name}. " \
      f"color: {color}. brand: {brand}."

text_embedding = text_model.encode(doc, normalize_embeddings=True)  # 384-d
\end{lstlisting}

Each photograph is embedded independently and stored on its own row, rather than
averaged into a single per-item vector. Averaging would blur a decisive
close-up together with an out-of-focus wide shot; keeping them separate lets the
scoring stage take the \emph{maximum} similarity over the cross product of the
two items' photographs --- best shot against best shot. A single matching angle
is strong evidence, and off-angle photographs should not be allowed to dilute
it. An item with no photograph simply has no image vector, a case the fusion
step handles explicitly.

Both models are loaded once per worker process as singletons and warmed at
startup, so their load cost is paid on deployment rather than on the first
report of the day.

\subsection{Retrieving candidates}

The engine never scores the whole table. Retrieval is a funnel with a cheap
relational stage followed by an approximate nearest neighbour stage.

\paragraph{Stage A --- relational pre-filter.} Candidates must be of the
opposite \code{type}, still \code{open}, and carry a \code{lost\_or\_found\_at}
date inside a plausible window. The window is deliberately wide and symmetric
(30 days by default): a found report can \emph{precede} the corresponding lost
report, because people frequently pick something up before its owner notices
the loss.

\paragraph{Stage B --- approximate nearest neighbour search.} Over the
pre-filtered set, \code{pgvector}'s cosine operator \code{<=>} retrieves the top
$K$ candidates using the HNSW index.

\begin{lstlisting}[language=SQL,caption={Candidate retrieval, filters and vector search in one query},label={lst:ann}]
SELECT i.id,
       1 - (i.text_embedding <=> :query_vec) AS vector_sim
FROM items i
WHERE i.type   = :opposite_type
  AND i.status = 'open'
  AND i.lost_or_found_at BETWEEN :date_lo AND :date_hi
  AND i.text_embedding IS NOT NULL
ORDER BY i.text_embedding <=> :query_vec    -- uses the HNSW index
LIMIT :k;                                    -- k = 50
\end{lstlisting}

\paragraph{The second route.} A text search alone systematically loses one case:
a good photograph attached to a two-word description. Such an item is never
retrieved by the text query, so it can never be scored, so it can never be
suggested --- no matter how obviously the photographs match. The engine
therefore runs a \emph{second, independent} retrieval over
\code{item\_images.image\_embedding}, and adds any item found that way to the
candidate set, scoring it on both modalities afterwards. The two routes are
unioned, which means an item can surface through strong wording \emph{or}
through a strong photograph, and is then judged on everything available.

\subsection{Calibrating the image similarity}
\label{sec:calibration}

This was the most instructive finding of the implementation, and it is worth
reporting because the naive version fails silently rather than loudly.

CLIP image embeddings are \emph{anisotropic}: they do not spread over the unit
sphere but occupy a comparatively narrow cone. This is a known property of
representations learned by contrastive and self-supervised
objectives~\cite{ethayarajh2019anisotropy}, and in multimodal contrastive models
it is compounded by a systematic separation between the image and text
manifolds~\cite{liang2022modalitygap}. The practical consequence,
measured on the platform's own corpus, is that two \textbf{completely unrelated}
photographs still score a cosine similarity of roughly $0.65$ to $0.71$, while
the same object photographed twice scores around $0.96$. Feeding the raw cosine
into a weighted average therefore treats an irrelevant photograph as roughly
$70\%$ agreement, which was enough to lift weak text matches over the
persistence threshold and put unrelated objects in front of users.

The fix is to rescale the interval $[f, 1]$ onto $[0, 1]$, where $f$ is a
configured floor set to the observed unrelated-pair level:

\begin{equation}
s_{\text{img}} = \operatorname{clamp}_{[0,1]}\!\left(\frac{s^{\text{raw}}_{\text{img}} - f}{1 - f}\right),
\qquad f = 0.65 .
\label{eq:calibrate}
\end{equation}

Under this mapping a raw $0.69$ becomes $0.11$ and a raw $0.96$ becomes $0.89$
--- the discrimination that the raw numbers were hiding. Clamping at zero
matters as much as the rescaling: it means \dq{less alike than two random
photographs} reads as \emph{no evidence} rather than as negative evidence
capable of dragging down an otherwise sound textual match.

\subsection{Fusing the scores}

Two adjustments are applied before fusion.

\paragraph{Lexical blending.} Embeddings capture meaning but blur exact tokens.
\dq{iPhone 14 Pro} and \dq{iPhone 13 Pro} sit almost on top of each other in
vector space, yet that difference is the entire identification. PostgreSQL's
full-text search scores them apart, so the semantic similarity is blended with a
lexical score, weighted towards the vector because exact overlap is precise when
present but absent for most genuine matches --- two people describing one object
in different words is the reason embeddings are in the system at all.

\begin{equation}
s_{\text{text}} = \operatorname{clamp}_{[0,1]}\!\big((1 - w_{\ell})\, s_{\text{vec}} + w_{\ell}\, s_{\text{lex}}\big),
\qquad w_{\ell} = 0.25 .
\label{eq:lexical}
\end{equation}

\paragraph{Weighted fusion with redistribution.} A missing photograph must not be
scored as a \emph{dissimilar} photograph. Renormalising by the weights that
actually fired means an item without a photograph is judged purely on its text,
at full scale, instead of being capped at $w_{\text{text}} = 0.5$ and ranking
below every photographed item however well it matches:

\begin{equation}
s_{\text{comb}} =
\frac{w_{\text{text}}\, s_{\text{text}}\, \mathbf{1}_{\text{text}}
    + w_{\text{img}}\, s_{\text{img}}\, \mathbf{1}_{\text{img}}}
     {w_{\text{text}}\, \mathbf{1}_{\text{text}}
    + w_{\text{img}}\, \mathbf{1}_{\text{img}}} ,
\label{eq:fusion}
\end{equation}

\noindent where $\mathbf{1}_{m}$ is the indicator that modality $m$ is present
for both items of the pair.

\subsection{Computing confidence and explanations}

The combined score measures embedding proximity. \emph{Confidence} is the
user-facing estimate that the pair really is the same object, and it starts from
the combined score and adds small bounded rewards for independent
corroboration:

\begin{equation}
c = \operatorname{clamp}_{[0,1]}\!\Big(s_{\text{comb}}
  + b_{\text{cat}} + b_{\text{col}} + b_{\text{brand}} + b_{\text{wil}}
  + b_{\text{time}} \cdot \tau + b_{\text{margin}} \cdot \mu \Big).
\label{eq:confidence}
\end{equation}

The boosts are deliberately small (Table~\ref{tab:params}). Together they can
lift a plausible match above the notification threshold, but they can never
manufacture a match out of an unrelated pair --- which is the failure that
matters more. A false \dq{we found your wallet} costs far more trust than a
missed suggestion costs opportunity.

Several details of the implementation follow from that asymmetry:

\begin{itemize}[leftmargin=1.4em,itemsep=3pt]
  \item \textbf{Category agreement counts only when both sides declared one.}
        Agreement on \emph{unknown} is not agreement, and a measurable share of
        real reports carry no category at all.
  \item \textbf{Proximity is measured by wilaya, not by radius.} The report form
        does not ask for coordinates and no report in the corpus carries any, so
        a haversine term would have been dead code. Wilayas are stored as
        numeric codes rather than names, precisely so that a location can be
        rendered in Arabic, French or English without touching stored data.
        \emph{Known limitation:} the implementation currently carries the 58
        wilayas in force since December 2019, whereas eleven further wilayas
        were created on 16~November 2025, bringing the national total to
        69~\cite{wilayas2025}. Because the codes are official administrative
        numbers and the new ones extend the sequence rather than renumbering it,
        updating the list is an additive change to one module and its client
        mirror; it does not affect stored data or the matching logic.
  \item \textbf{Temporal plausibility is directional.} An object is found
        \emph{after} it is lost, so the reverse ordering is penalised --- but
        halved rather than zeroed, because people notice a loss late and
        remembered dates are approximate.
  \item \textbf{Distinctiveness is rewarded.} A top candidate standing clearly
        above the runner-up is more trustworthy than the same score in a
        five-way tie, where the similarity is probably measuring \dq{generic
        black wallet} rather than \emph{this} wallet. Because this term depends
        on the rest of the field, it can only be applied once every candidate
        has been scored.
\end{itemize}

\paragraph{Explanations as reason codes.} The interface shows the confidence as
a percentage together with the reasons that produced it. Those reasons are
emitted by the worker as \textbf{codes with parameters} --- \code{same\_category},
\code{same\_wilaya}, \code{time\_close} with a day count --- and never as
English prose. The worker cannot know which of three languages the eventual
reader uses, and this is the product's flagship surface: an Arabic-speaking user
reading \dq{found 1 day after lost} in English is a visible failure. The
frontend translates the codes through the same message catalogues that drive the
rest of the interface.

\begin{table}[htbp]
\centering
\caption{Matching parameters, all environment-configurable}
\label{tab:params}
\small
\begin{tabularx}{\textwidth}{@{}l l X@{}}
\toprule
\textbf{Parameter} & \textbf{Value} & \textbf{Meaning} \\
\midrule
\code{MATCH\_TOPK}            & 50   & Candidates retrieved per modality before scoring \\
\code{MATCH\_DATE\_SLACK\_DAYS} & 30 & Half-width of the plausible date window \\
\code{MATCH\_W\_TEXT}         & 0.50 & Weight of the textual modality \\
\code{MATCH\_W\_IMAGE}        & 0.50 & Weight of the visual modality \\
\code{MATCH\_W\_LEXICAL}      & 0.25 & Share of full-text score blended into the text similarity \\
\code{MATCH\_IMAGE\_SIM\_FLOOR} & 0.65 & Raw CLIP cosine treated as zero evidence (Eq.~\ref{eq:calibrate}) \\
\code{MATCH\_BOOST\_CATEGORY} & 0.05 & Both sides declare the same category \\
\code{MATCH\_BOOST\_COLOR}    & 0.04 & Colours agree \\
\code{MATCH\_BOOST\_BRAND}    & 0.04 & Brands agree \\
\code{MATCH\_BOOST\_WILAYA}   & 0.05 & Same wilaya \\
\code{MATCH\_BOOST\_TIME}     & 0.04 & Scaled by temporal plausibility \\
\code{MATCH\_BOOST\_MARGIN}   & 0.05 & Scaled by the gap to the runner-up \\
\code{CONF\_PERSIST}          & 0.55 & Below this, the match is not stored at all \\
\code{CONF\_NOTIFY}           & 0.70 & At or above this, both owners are notified \\
\code{CONF\_STRONG}           & 0.85 & Interface shows a \dq{high confidence} badge \\
\bottomrule
\end{tabularx}
\end{table}

Keeping every one of these values in configuration rather than in code is not
tidiness. These are exactly the parameters a future model will learn from
accumulated feedback --- by fitting the weights on confirmed and rejected pairs
and mapping the resulting scores onto true probabilities through Platt scaling
or isotonic regression~\cite{zadrozny2002calibration,guo2017calibration} --- and
holding them behind a single scoring interface means the learned version can
replace the hand-set version without touching the API, the database or the
interface.

\subsection{Persistence, retraction and notification}

Three behaviours make the engine safe to re-run, which matters because matching
is triggered again whenever either side of a pair is edited.

\begin{itemize}[leftmargin=1.4em,itemsep=3pt]
  \item \textbf{Idempotence.} The \code{matches} table is unique on the (lost,
        found) pair, so a re-run \emph{upserts} updated scores rather than
        accumulating duplicates.
  \item \textbf{Retraction.} Any pair that this pass no longer supports is
        expired. Without this, re-scoring could only ever add suggestions; an
        item whose entire candidate pool had closed would go on displaying
        matches it no longer has.
  \item \textbf{Notification on transition only.} A user is notified when a pair
        \emph{crosses} the notification threshold, not on every upsert.
        Otherwise correcting a typo in a description would send a fresh
        \dq{we may have found your item} for every pair, every time.
\end{itemize}

\subsection{End-to-end sequence}

Figure~\ref{fig:sequence} traces one complete cycle, from the moment a report is
filed to the moment contact details are released. Two things are worth reading
off it. The reporter receives \code{201~Created} long before any model has run
--- the response does not wait on inference. And the final disclosure is
triggered by a human decision, not by the score: the worker's notification ends
at \dq{a claim is awaiting your review}, and nothing further happens until the
reporter approves.

\begin{figure}[htbp]
\centering
\begin{tikzpicture}[x=1cm,y=1cm]

  % --- participants ---
  \node[actorbox] (a) at (-6,0) {Reporter};
  \node[actorbox] (b) at (-3,0) {API};
  \node[actorbox] (c) at ( 0,0) {Worker};
  \node[actorbox] (d) at ( 3,0) {Database};
  \node[actorbox] (e) at ( 6,0) {Counterpart};

  % --- lifelines ---
  \draw[lifeline] (-6,-0.35) -- (-6,-10.1);
  \draw[lifeline] (-3,-0.35) -- (-3,-10.1);
  \draw[lifeline] ( 0,-0.35) -- ( 0,-10.1);
  \draw[lifeline] ( 3,-0.35) -- ( 3,-10.1);
  \draw[lifeline] ( 6,-0.35) -- ( 6,-10.1);

  % --- messages ---
  \draw[msg] (-6,-1.0) -- (-3,-1.0) node[msglabel,midway] {POST /items + photos};
  \draw[msg] (-3,-1.6) -- ( 3,-1.6) node[msglabel,midway] {insert item, images, storage keys};
  \draw[msg] (-3,-2.2) -- ( 0,-2.2) node[msglabel,midway] {enqueue \dcode{embed\_item}};
  \draw[retmsg] (-3,-2.8) -- (-6,-2.8) node[msglabel,midway] {201 Created};

  \draw[msg] ( 0,-3.6) -- ( 3,-3.6) node[msglabel,midway] {write 384-d and 512-d vectors};
  \draw[msg] ( 0,-4.2) -- ( 3,-4.2) node[msglabel,midway] {ANN + full-text over opposite pool};
  \draw[retmsg] ( 3,-4.8) -- ( 0,-4.8) node[msglabel,midway] {top-$K$ candidates};

  \draw[msg] (0,-5.35) -- ++(0.9,0) -- ++(0,-0.5) -- ++(-0.9,0);
  \node[font=\tiny, anchor=west] at (1.05,-5.6) {calibrate, fuse, score, apply margin};

  \draw[msg] ( 0,-6.5) -- ( 3,-6.5) node[msglabel,midway] {upsert \dcode{matches}};
  \draw[msg] ( 0,-7.1) -- (-6,-7.1) node[msglabel,midway] {notification: possible match};
  \draw[msg] ( 0,-7.7) -- ( 6,-7.7) node[msglabel,midway] {notification: possible match};

  \draw[msg] ( 6,-8.4) -- (-3,-8.4) node[msglabel,midway] {submit claim + answers to verification questions};
  \draw[msg] (-3,-9.0) -- (-6,-9.0) node[msglabel,midway] {claim awaiting review};
  \draw[msg] (-6,-9.6) -- (-3,-9.6) node[msglabel,midway] {approve};
  \draw[msg] (-3,-10.0) -- ( 6,-10.0) node[msglabel,midway] {contact details released};

\end{tikzpicture}
\caption{Sequence from report to handover. Inference is entirely off the request
path, and no identity is disclosed without the reporter's explicit approval.}
\label{fig:sequence}
\end{figure}

\section{Presentation and implementation of the user interfaces}

\towrite{Capture the screenshots listed in \code{thesis/figures/README.md} and
drop them into \code{thesis/figures/}. Each figure below switches from a grey
placeholder to the real image automatically once the file exists --- no LaTeX
edit needed. Keep the browser at roughly 1440\,px wide, use the seeded demo
data so the grids are full, and take at least one capture in Arabic to
evidence the right-to-left layout.}

\subsection{Functional coverage}

Figure~\ref{fig:usecase} summarises what the implemented system lets each actor
do. It is a map of what exists, not of what was planned: administration and
moderation surfaces are specified but not yet shipped, and are therefore absent
from the diagram. A partner institution uses the same reporting path as an
individual finder --- publishing a register is not a separate feature but a
volume of ordinary found reports --- which is why no dedicated institutional
use case appears.

\begin{figure}[htbp]
\centering
\begin{tikzpicture}[x=1cm,y=1cm]

  \node[systemframe, minimum width=7.2cm, minimum height=8.0cm] (sys) at (0,2.8) {};
  \node[font=\scriptsize\bfseries\color{uamobblue}, anchor=north] at (0,6.65) {\projectname{} platform};

  \node[usecase, minimum width=34mm] (u1) at (0,6.0) {Browse and search reports};
  \node[usecase, minimum width=34mm] (u2) at (0,5.1) {Register and authenticate};
  \node[usecase, minimum width=34mm] (u3) at (0,4.2) {Report a lost or found item};
  \node[usecase, minimum width=34mm] (u4) at (0,3.3) {Upload photographs};
  \node[usecase, minimum width=34mm] (u5) at (0,2.4) {View ranked match suggestions};
  \node[usecase, minimum width=34mm] (u6) at (0,1.5) {Submit a claim with answers};
  \node[usecase, minimum width=34mm] (u7) at (0,0.6) {Approve or reject a claim};
  \node[usecase, minimum width=34mm] (u8) at (0,-0.3) {Receive notifications};

  \pic (av) at (-5.9,5.9) {actor};
  \node[font=\tiny, align=center] at (-5.9,4.3) {Visitor};

  \pic (au) at (-5.9,2.6) {actor};
  \node[font=\tiny, align=center] at (-5.9,1.0) {Registered\\user};

  \pic (ai) at (5.9,4.6) {actor};
  \node[font=\tiny, align=center] at (5.9,3.0) {Partner\\institution};

  \pic (aw) at (5.9,1.0) {actor};
  \node[font=\tiny, align=center] at (5.9,-0.6) {Matching\\worker};

  \draw[relline] (-5.9,5.4) -- (u1.west);
  \draw[relline] (-5.9,5.4) -- (u2.west);

  \draw[relline] (-5.9,2.1) -- (u3.west);
  \draw[relline] (-5.9,2.1) -- (u4.west);
  \draw[relline] (-5.9,2.1) -- (u5.west);
  \draw[relline] (-5.9,2.1) -- (u6.west);
  \draw[relline] (-5.9,2.1) -- (u7.west);
  \draw[relline] (-5.9,2.1) -- (u8.west);

  \draw[relline] (5.9,4.1) -- (u3.east);
  \draw[relline] (5.9,4.1) -- (u4.east);
  \draw[relline] (5.9,4.1) -- (u7.east);

  \draw[relline] (5.9,0.5) -- (u5.east);
  \draw[relline] (5.9,0.5) -- (u8.east);

\end{tikzpicture}
\caption{Use cases implemented in the prototype. The matching worker is a system
actor: the suggestions and notifications a user sees are produced without any
human initiating them.}
\label{fig:usecase}
\end{figure}

\subsection{Home page}

The landing page states the proposition in one line, offers the two entry points
that matter --- report something lost, report something found --- and shows the
most recent reports so that a first-time visitor sees an active platform rather
than an empty form. The language switcher and the light/dark toggle are
available before authentication.

\screenshot{home.png}{Public landing page with the two report entry points}{home}

\subsection{Registration and authentication}

Registration takes an e-mail, a full name and a password. Passwords are hashed
with Argon2, the memory-hard function standardised in RFC~9106~\cite{argon2rfc}.
Sessions use short-lived access tokens (15 minutes) paired with
rotating refresh tokens valid for 30 days, and every rotation invalidates its
predecessor.

Reuse of an already-rotated refresh token is treated as potential theft and
revokes the session chain --- but only outside a brief grace window. That window
exists because the benign case is common: a second browser tab, a page reload in
mid-refresh, or a retried request will legitimately replay the previous token,
and without the grace period a single stale replay would log the user out of
every device they own.

\screenshot{auth-login.png}{Sign-in, with client-side validation from the shared Zod schema}{auth}

\subsection{Reporting an item}

One guided form serves both directions, differing only in wording and in the
date label. It collects a title, a free-text description, a category, colour,
brand, wilaya, the date, up to two verification questions, and up to five
photographs.

Uploads are validated on content rather than on file extension --- declared MIME
type, magic bytes and size --- then re-encoded with Pillow to WebP at a longest
edge of 1600\,px. A typical phone photograph of three to five megabytes comes
down to roughly 150--250\,kilobytes, which is the difference between a few
hundred and a few thousand photographs on a free-tier storage bucket, and the
1600\,px bound still exceeds every size at which the interface renders an image.

\screenshot{report-form.png}{Reporting a found item: details, photographs and verification questions}{report}

\subsection{Browsing and searching}

The public browse views list lost and found reports with filters on category,
wilaya and date, and a keyword search served by the PostgreSQL full-text index.
Keyword search is a genuine fallback rather than a lesser version of the
matching engine: a user who knows the brand name written on the object wants
exactly that token, not its neighbourhood in vector space.

\screenshot{browse.png}{Browse view with category, wilaya and date filters}{browse}

\subsection{Item detail and match suggestions}

The detail page shows the report, its photographs, and --- for the owner --- the
ranked suggestions produced by the engine. Each suggestion carries its
confidence as a percentage, a badge when it exceeds the strong threshold, and
the translated list of reasons that produced it. Showing the reasons is a
deliberate product decision: a bare percentage invites either blind trust or
blind dismissal, whereas \dq{same category, same wilaya, photographs strongly
similar, found one day after the loss} lets the user apply judgement the system
does not have.

\screenshot{item-matches.png}{Match suggestions with confidence and generated explanations}{matches}

\subsection{The claim workflow}

A match is a hypothesis, not an authorisation, so a confirmed match never
reveals anyone's contact details. Recovery goes through a separate claim:

\begin{enumerate}[leftmargin=1.6em,itemsep=3pt]
  \item The reporter sets up to two verification questions when publishing
        (\dq{Which side is it, left or right?}, \dq{What is the lock-screen
        photograph?}).
  \item A user who believes the item is theirs submits a claim answering those
        questions, with an optional message. The questions are snapshotted onto
        the claim, so the answers stay meaningful even if the reporter later
        edits them.
  \item The reporter reviews and approves or rejects.
  \item \textbf{Contact details are exchanged only on approval}, and only
        between those two people.
\end{enumerate}

One symmetric mechanism covers both directions: on a found item the loser proves
ownership, and on a lost item the finder proves possession. A partial unique
index permits only one \emph{pending} claim per person per item, which stops a
claimant from flooding a reporter while still allowing a fresh attempt after a
rejection.

\screenshot{claim.png}{Submitting a claim by answering the reporter's verification questions}{claim}

\subsection{Notifications and personal dashboard}

Notifications are created in-app when a match crosses the notification
threshold, when a claim is submitted, and when a claim is resolved. Each carries
a deep link to the item or match that caused it. The channel column already
distinguishes in-app from e-mail, so an e-mail transport can be added without a
schema change.

The dashboard collects the user's own reports with their live state --- the ML
pipeline status while embedding is in progress, the number of suggestions, and
any pending claims awaiting their decision.

\screenshot{notifications.png}{Notification centre}{notifications}
\screenshot{dashboard.png}{Personal dashboard: own reports, suggestions and pending claims}{dashboard}

\subsection{Trilingual interface and right-to-left support}

Every route is prefixed by a locale segment, and English, French and Arabic each
have a complete message catalogue. Arabic is rendered right-to-left; because the
styling uses logical rather than physical properties, the layout mirrors without
a parallel stylesheet.

Localisation reaches further than the interface strings. Wilayas are stored as
numeric codes and rendered per locale. Match explanations travel from the worker
as reason codes and are translated on the client. And the text encoder itself is
multilingual, which is what allows a description written in Arabic to match a
description of the same object written in French --- the single most important
consequence of the model choice in Section~\ref{sec:ml-stack}.

\screenshot{arabic-rtl.png}{The same interface in Arabic, mirrored right-to-left}{rtl}

\section{Testing and validation}

Testing is concentrated where a defect would be either invisible or expensive.

\begin{itemize}[leftmargin=1.4em,itemsep=3pt]
  \item \textbf{Scoring is pure and unit-tested.} Fusion, calibration,
        confidence and the margin rule are plain functions over plain values ---
        no database session, no model, no PyTorch. The entire scoring policy can
        therefore be tested in milliseconds, without loading half a gigabyte of
        weights, and the tests can be imported from the API side as well.
  \item \textbf{Integration tests exercise the real HTTP surface} for
        authentication and the item lifecycle, through the same validation and
        authorisation dependencies that run in production.
  \item \textbf{Schema tests} pin the request and response contracts that the
        typed frontend client is generated against.
  \item \textbf{Security tests} cover token issuance, expiry, rotation and the
        reuse-detection path, including the grace window.
  \item \textbf{Health endpoints} separate liveness from readiness: readiness
        checks the database and, when configured, Redis.
\end{itemize}

The most valuable validation, however, is not automated. Working through the
seeded demonstration corpus by hand is what exposed the CLIP anisotropy
described in Section~\ref{sec:calibration} --- unrelated photographs were being credited with
roughly $70\%$ agreement, and no unit test would have flagged that, because the
code was doing exactly what it had been written to do. The number was wrong, not
the arithmetic.

\section{Deployment and packaging}

The whole system runs from a single \code{docker compose up}. Five services are
defined --- \code{postgres} (with \code{pgvector}), \code{redis}, \code{api},
\code{worker} and \code{frontend} --- with named volumes for database files,
Redis persistence, uploaded media, and the Hugging Face model cache. Caching the
model weights in a volume matters more than it appears to: without it, every
container rebuild re-downloads several hundred megabytes of weights before the
worker can start.

The API and the worker are built from the same Dockerfile but install different
dependency sets. The API never imports PyTorch, so the ML requirements are
confined to the worker stage; and PyTorch is installed from the CPU wheel index,
because the default distribution bundles CUDA and pulls roughly 2.5\,GB of NVIDIA
libraries that this deployment cannot use, against roughly 200\,MB for the CPU
build.

Two configuration choices keep low-cost hosting viable. Setting
\code{REDIS\_URL} to an empty string runs the API queue-free, with readiness
reporting Redis as skipped rather than failing --- enough to deploy the browsing
and reporting surface as a single service on a free tier. And the
\code{StorageBackend} abstraction means moving from local disk to S3-compatible
object storage is a configuration change, not a code change.

\section{Conclusion}

This chapter has described the prototype as it stands: a modular-monolith
backend with all inference isolated on a background worker, a PostgreSQL
database doing relational filtering, full-text search and vector similarity in
one engine, and a trilingual Next.js interface covering the full cycle from
report to recovery.

The engineering that distinguishes the result is concentrated in the matching
engine, and specifically in the parts that keep it honest --- calibrating a raw
CLIP cosine that was silently inflating every score, redistributing fusion
weights so that a missing photograph is not read as a dissimilar one, blending a
lexical term so that \dq{iPhone 14} and \dq{iPhone 13} stop being
indistinguishable, retracting suggestions the evidence no longer supports, and
keeping contact details behind a verification handshake that a similarity score
alone can never open. Every scoring constant remains in configuration, and every
confirmed or rejected suggestion is persisted as a labelled example, so the
hand-tuned policy described here is positioned to be replaced by a learned and
calibrated one without disturbing anything around it.
